数据中心内的内存事务面对另一组成本：对象已经在DRAM中，远程读取一个小对象，却可能先经过远端CPU调度、消息分发和协议处理。FaRM提供跨机器的全局地址空间，让应用在事务中读写本地或远端对象，并用RDMA缩短通信路径。数据的放置、并发验证、复制提交和故障恢复围绕这条路径共同设计~\cite{farm2014,farm2015}。

\subsection{对象、region与单边访问}

FaRM把内存组织为region。一个region有稳定标识、primary及备份，对象地址由region标识和其中的偏移组成。运行时根据当前映射找到primary：对象在本机就用本地内存访问，在远端就取得相应RDMA引用。对象读从primary取得，备份则承接复制与故障后的保存责任。

远程直接内存访问RDMA使网卡按已注册的内存区域和权限读写远端内存。内存注册建立网卡需要的映射，并固定相关页面。单边读取由远端网卡完成，远端应用线程不用先收包再复制对象；调用方通过完成队列观察操作完成。事务代码、对象更新和日志处理仍由相应CPU执行。

\CNineteenFigure{ch19-17}{region映射把对象地址连接到primary，单边访问和日志软件处理分别发生}{fig:19-farm-address}

FaRM也用单边写构造FIFO消息队列和事务日志。每个发送者与接收者之间有对应环形缓冲区，位于接收端内存。发送者向尾部写记录，接收线程轮询处理，再逐步通知可复用位置。网卡已经确认日志写入时，接收线程可能尚未解释或应用该事务记录。

注册内存、页表及队列对也占用网卡资源。大页、连接共享和相关对象的局部放置降低这些成本。把常在同一事务访问的对象放到同一primary及同一组备份上，能使更多锁定、验证和读写留在本地。分散全部对象会增加网络交互，过度集中又会形成热点，放置策略据真实访问集合安排。

region的位置映射可以用不同结构维护。共享远程内存设计按多条虚拟环的一致性哈希选择primary和副本，成员信息使各机器能够本地计算位置。集成事务复制设计由配置管理器分配region并选择副本，显式考虑容量、局部性与独立故障域，映射随region复制。集中选择增加了放置灵活性，配置管理器同时承担分配与变更的控制责任~\cite{farm2014,farm2015}。

\subsection{一次对象读取怎样排除混合字段}

RDMA读取可能跨多个缓存行，并与CPU更新交错。只检查对象头的一个版本号，可能取得旧头和新数据。FaRM在对象头与后续缓存行安排版本信息，读取以后核对头部锁状态及各缓存行版本。成功检查使返回内容对应一个完整对象状态，失败则丢弃本次对象并重读~\cite{farm2014}。

\CNineteenFigure{ch19-18}{头部及缓存行版本识别混合对象，incarnation另外识别地址释放后的复用}{fig:19-farm-version}

例如，一个对象记录收款C1及金额3000，更新改成另一完整状态。读取者得到旧C1和新金额时，各字段都能解析，却没有对应同一对象版本。整对象验证把这种组合挡在业务计算之前。写入方更新缓存行、版本与头部的先后，也通过内存屏障约束，使读取方的检查有明确依据。

地址复用则需要额外信息。地址$a$原先保存X，X释放后又保存Y，迟到指针继续指向$a$。incarnation表示对象化身，指针保存预期化身，读取后比较以识别释放与复用。共享远程内存设计使用包含地址、大小和预期化身的指针，并在释放对象时推进化身。不同大小对象的复用还要等待相关旧访问退出，避免对象头已被其他布局覆盖。

这些规则依赖缓存一致DMA、写入次序、版本位数及读取时间界。较短缓存行版本可能回绕，协议以更新速度和最大读取时间的界排除一次读取跨越同版本再现。地址负责定位，版本负责并发，化身负责复用，业务$S$负责关联一次购买，各自保留自己的比较对象。

\subsection{乐观执行与提交验证}

乐观并发控制OCC把事务的大量计算放在取得全部锁以前执行。它适合冲突较少、事务较短的访问：多数尝试先读取并计算，临近提交时再检查依据是否变化。冲突频繁时，已经完成的计算可能被反复丢弃，重试成本及热点写锁就成为重要约束。

执行阶段读取对象，记录地址和版本，形成读集；准备修改的对象进入写集，新值暂存在本地事务上下文。写集通常包含地址、预期版本及新值。此时尚未改变共享对象。读取单个对象可以完整，但读取多个对象期间其他事务仍会提交，执行代码因此要能安全处理暂时不一致的跨对象观察~\cite{farm2015}。

\CNineteenFigure{ch19-35-occ}{执行阶段保存读版本与本地写集，提交时锁定写对象并验证只读依据}{fig:19-occ}

设事务读取报价Q7的版本8、库存A的版本12和库存B的版本20，准备各预留一件。报价只参与判断，没有进入写集。若另一事务修改报价，原事务即使能锁住A和B，也要在读集验证中发现报价版本变化，放弃原计算。只验证待写对象，会遗漏决定购买条件的只读依据。

FaRM提交时按原版本锁定写对象。对象已经被锁或版本改变，就abort；全部写锁成功以后，再验证读而未写的对象版本。验证失败时释放相应锁并放弃本次暂存写入。允许重试的事务重新读取和计算，应用则沿原业务绑定保存$S$及原输入。

另取一条初始库存为一件、版本12的历史。两个事务都读到这一版本，第一个锁定并提交后，库存为零、版本推进。第二个按版本12锁定时失败，旧读取不能覆盖已经完成的预留。若第二次到来的是同一$S$的网络重试，业务仲裁先找到原记录，回放原结果或继续原责任。内部OCC尝试可以多次，业务购买身份保持同一项。

\subsection{复制日志怎样支撑提交}

并发验证解释事务能否成立，复制记录解释提交结果怎样跨故障保存。FaRM的共享远程内存结构把prepare发送到写对象的primary和副本：primary锁定对象，primary与副本记录消息后答复。协调者再验证读集，成功后先向副本发送commit，再向primary发送commit，由primary更新对象并解锁。副本日志保存于SSD，以少量非易失内存缓冲日志写入~\cite{farm2014}。

\CNineteenFigure{ch19-34-ssd-commit}{复制日志落SSD的结构先保存prepare，再验证读集，并先提交副本后更新primary}{fig:19-farm-ssd}

若把整个数据与日志都放在可恢复的非易失DRAM中，就能进一步把事务提交与主备复制合并。FaRM的集成结构让执行事务的应用线程所在机器承担协调者，协调者自身没有逐事务Paxos状态机。它直接在primary与备份之间保存足够记录，让正常提交和故障恢复从这些记录判断结果~\cite{farm2015}。

第一步Lock：协调者向各写对象primary的日志写LOCK，包含事务身份、所修改region集合、地址、原版本和新值。primary按预期版本尝试取得写锁，并返回结果。这条路径在锁阶段只联系primary，备份在验证成功后的提交阶段接收记录。

第二步Validate：协调者读取只读对象的版本，确认执行依据没有改变。小读集用单边RDMA验证，某个primary上的待验证对象较多时，可以用RPC合并验证，减少大量分散读取。锁定失败或验证失败，按协议abort。

第三步Commit backups：协调者把COMMIT-BACKUP写到每个相关备份的非易失日志，并等待全部网卡确认。记录携带写对象的新值及事务信息，备份CPU可以稍后批量处理。

第四步Commit primaries：所有备份写入都确认以后，再向各primary日志写COMMIT-PRIMARY。至少一份这样的网卡确认到达，或者相应本地日志写完成以后，协调者才向应用报告提交完成。primary处理记录时更新对象、推进版本并解锁，使新访问看见提交结果。

\CNineteenFigure{ch19-19}{集成提交先锁定和验证，再保存全部备份提交记录，最后推进primary}{fig:19-farm-commit}

先等待全部备份，保护了新值的存活条件。若某region的一个备份还没有提交记录，primary已经暴露修改，随后primary、协调者和其他相关副本一同失败，存活的这份备份就可能只有旧值。按上述顺序，primary暴露新值以前，相应提交内容已经到达全部要求的备份。

向应用报告成功前再等待至少一份primary提交记录，则使验证成功留下足够证据。读集主要保存在协调者，只有LOCK存活时，恢复者未必能证明验证已经通过。成功时点连到提交记录的保存位置，使后续故障恢复能够维护已经对外完成的结果。

最后延后截断。协调者等各primary确认，再在后续记录中携带已完成事务身份，推进日志截断，备份在截断时把更新应用到对象副本。单边写没有逐次等待备份CPU分配空间，协调者便在开始提交前预留提交及截断记录所需的日志空间。日志满时，释放空间的记录也有位置可写。

\subsection{非易失保存与故障域}

这里的非易失DRAM建立在明确硬件路径上。FaRM采用本地电池继续供电，在掉电后把内存转存到SSD，恢复供电以后再读回。正常事务路径从而可以把写入多个非易失内存副本作为保存条件，免去每次提交同步写SSD的等待~\cite{farm2015}。

\CNineteenFigure{ch19-36-persistence}{正常事务写入非易失DRAM，整机掉电时由电池支持转存SSD，再恢复内存与日志}{fig:19-farm-persistence}

持久性依赖电池能量、转存程序、SSD容量及副本独立故障域。普通易失DRAM里的RDMA写确认，只说明相应内存写入完成。要采用论文的掉电恢复结论，数据与日志就要处于该非易失保存基础上。

协议按每个对象一个primary和$f$个备份组织保存，在相应故障条件下处理至多$f$份副本失去非易失内容。完整集群掉电的恢复又依赖保存内容能够读回。继续提供服务还要求：存活分区包含多数机器，能连接ZooKeeper多数，并且每个对象至少有一份副本。进程失效停止、时钟漂移有界及最终通信延迟有界，分别进入安全与进展的假设。

这些条件解释了它与跨地域复制数据库的取舍。集成主备事务减少常见路径的复制消息和远端CPU参与，配置协议却需要精确成员关系，硬件也要承担非易失保存与RDMA访问条件。跨地域Paxos复制把各组的复制决定与跨组事务分开，广域往返和时间等待进入相应路径。选择由数据局部性、故障域、硬件及访问持续时间共同决定。

\subsection{配置交接与事务恢复}

FaRM用配置管理器和ZooKeeper共同保存配置变化。配置包含单调配置身份、成员、故障域和管理者。ZooKeeper参与配置交接的共同选择，普通事务直接访问对象及事务日志。新配置提升受影响region的存活备份，并重新安排副本。

单边访问使成员变化进入数据路径的正确性。远端CPU不能在每次RDMA到达时检查租约，网卡可能已经接受旧请求。FaRM通过精确成员关系，使新配置中的机器在允许对象修改前同意成员集合，停止向已移除成员发新操作，并忽略来自被移除成员的不合条件完成。配置交接还等待相应旧租约边界~\cite{farm2015}。

\CNineteenFigure{ch19-20}{新配置先排空旧日志、恢复锁与事务状态，再逐步恢复访问和数据冗余}{fig:19-farm-recovery}

旧配置记录又不能一概丢弃。网卡确认可能已经使协调者回复成功。新配置排空相关日志，记下LastDrained边界，收集受影响事务及其region集合，再补回所需记录并恢复写锁。受影响region在锁恢复前暂停新的对象引用，防止新事务越过尚未确定的更新。

恢复者根据各region保存的记录投票。存在COMMIT-PRIMARY等提交证据，必须维护提交；仅有LOCK则表明锁定准备，并没有单独证明读集验证成功。没有记录时，还要区分日志已经合法截断和这份事务从未到达。把合法截断和从未到达都当作尚未提交，会遗漏已经形成并完成截断的提交事实。

\CNineteenFigure{ch19-37-recovery-votes}{恢复收集各region的存活日志，以提交证据、准备及截断状态恢复同一事务结果}{fig:19-farm-votes}

恢复协调者收到任一region的commit-primary票就决定commit。否则等待全部相关region：至少一份commit-backup，其他票均为lock、commit-backup或truncated时，按协议commit；其余组合决定abort。决定发到参与副本，再等待确认并截断恢复记录。恢复责任由配置和事务身份分配，能够跨机器与线程并行推进~\cite{farm2015}。

锁恢复完成后，region可以恢复受保护的访问，尚未决定的对象继续被对应锁约束。大量数据补副本与分配器恢复随后并行进行，前台无须等全部字节搬运完才继续。恢复流量和正常事务仍共享内存、网卡与CPU，系统为锁恢复和数据复制分别安排预算。

\subsection{把完整购买放进事务对象}

若用这类平台保存购买，完整$S$映射到原输入、实际P/C/T和原结果，必要库存进入相应读写集合。OCC验证版本与锁，复制协议保存提交证据，业务代码仲裁同一$S$的输入是否相同。只把库存计数放入事务，无法为网络重试恢复原购买的受理关联。

执行阶段的本地计算和暂存写入可以随冲突放弃，平台已经记账的付款则独立存在。应用先提交原C/T及发送资格，由可替换工作者携带原scope和\code{pay-C}付款。结果未知时查询原操作，再把确定结果写回购买权威。

较大的图像、分布式事务和外部付款在这套结构中分别有明确保存位置。批量对账读取固定manifest或数据库时间，任务运行时从持久责任取得原身份，执行者退出以后，后继者沿同一记录继续。
\FloatBarrier
